% Capítulo 1: Introducción, motivación y objetivos.

Este primer capítulo sitúa el problema de la selección de
características en el contexto de la Ciencia de Datos y delimita el
propósito del trabajo. En primer lugar se expone la motivación, que
conecta la reducción de dimensionalidad con la estabilidad, la
interpretabilidad y la validez de los modelos supervisados. A
continuación se fijan el objetivo general y los objetivos específicos
que guían el estudio.

\section{Motivación}

La selección de características es una de las etapas centrales del
aprendizaje automático supervisado. Su objetivo es identificar el
subconjunto de variables que conserva la capacidad predictiva del
problema reduciendo dimensionalidad, redundancia y coste
computacional~\cite{guyon2003}. En espacios de alta dimensionalidad,
una buena selección no es solo una optimización: condiciona la
estabilidad de los modelos, su interpretabilidad y la propia validez
de las conclusiones, como puso de manifiesto el desafío de selección
de características de NIPS~2003, un banco de referencia clásico para
problemas de alta dimensionalidad~\cite{guyon2004}.

El problema de seleccionar el mejor subconjunto de $k$ variables entre
$p$ candidatas es combinatorio y, en general, intratable de forma
exacta. Los métodos clásicos lo abordan mediante heurísticas: filtros
univariantes, métodos basados en regularización, importancias de
modelos de conjunto o criterios de relevancia--redundancia como
mRMR (\textit{minimum Redundancy Maximum Relevance})~\cite{peng2005}. En paralelo, la computación cuántica ofrece una
vía alternativa: formular la selección como un problema de
optimización binaria cuadrática sin restricciones (QUBO,
\textit{Quadratic Unconstrained Binary Optimization}) que combine
relevancia y redundancia en una única función objetivo, y resolverlo
sobre hardware cuántico~\cite{mucke2023}. En particular, las plataformas de átomos
neutros permiten codificar este tipo de problemas aprovechando el
fenómeno del bloqueo de Rydberg, que impide la excitación simultánea
de átomos vecinos y actúa de forma natural como una restricción de
exclusión entre variables redundantes~\cite{henriet2020}.

Este Trabajo de Fin de Grado analiza la viabilidad de un método
cuántico de selección de características (QFS, \textit{Quantum Feature
Selection}) basado en dicho fenómeno. El atractivo de esta vía es
computacional: QFS asigna un átomo (un cúbit) a cada variable
candidata y resuelve el compromiso relevancia--redundancia de forma
global sobre el hardware, de modo que un dispositivo de $n$ átomos puede
seleccionar características sobre un espacio de hasta $n$ variables en un
tiempo de ejecución muy corto, frente al coste combinatorio de la
búsqueda exacta clásica. Evaluar esa viabilidad de forma interpretable
exige caracterizar primero cada conjunto de datos: dónde reside la señal
predictiva, cómo se distribuye la redundancia, qué dimensionalidad
efectiva tiene y si las particiones permiten atribuir los resultados al
método y no al azar o a fugas de información. Con ese fin se construye
una referencia clásica que aporta, además de un baseline de rendimiento,
un sistema de coordenadas de relevancia, redundancia, estabilidad y coste
sobre el que situar el método cuántico.

De ahí la tesis que guía la memoria: el bloque clásico no es un mero
punto de comparación, sino el instrumento que permite leer el método
cuántico. Una caracterización clásica rigurosa anticipa qué componente
de QFS puede fallar en cada régimen de datos; y el contraste con el
óptimo exacto del problema de optimización subyacente permite, cuando
aparece un deterioro, separar si procede del criterio de información
mutua o del optimizador analógico que intenta alcanzarlo.

\section{Objetivos}

El objetivo general del TFG es analizar la viabilidad de un método de
selección de características que hace uso del bloqueo de Rydberg,
presente en tecnologías de átomos neutros, dentro de un protocolo
supervisado de Ciencia de Datos. La pregunta no es solo si QFS mejora
una métrica final, sino si puede competir con referencias clásicas
contrastadas y si sus resultados pueden atribuirse al criterio de
selección o al proceso de optimización. La propuesta del trabajo lo
desarrolla en seis objetivos específicos:

\begin{enumerate}
\item Revisar métodos clásicos de selección de características
  ampliamente utilizados y representativos de distintas familias de
  selección.
\item Seleccionar conjuntos de datos de referencia, diversos en número
  de muestras, dimensionalidad y tipo de variables, sobre los que
  evaluar los métodos de selección.
\item Realizar un estudio comparativo del rendimiento de los métodos
  clásicos sobre los conjuntos seleccionados.
\item Implementar y evaluar un método cuántico de selección de
  características basado en el bloqueo de Rydberg, conservando el núcleo
  de la formulación de referencia e incorporando adaptaciones operativas
  propias: un parámetro de separación $\beta$, una envolvente híbrida de
  preselección y la extensión del banco de pruebas a problemas
  multiclase y de alta dimensionalidad.
\item Comparar el rendimiento del método cuántico frente a los mejores
  selectores clásicos identificados, bajo un protocolo de evaluación
  común.
\item Analizar, mediante el óptimo exacto de la formulación QUBO, si las
  diferencias observadas entre el método cuántico y las referencias
  clásicas proceden del criterio de selección o del proceso de
  optimización analógica.
\end{enumerate}

El estudio clásico (objetivos 1--3) no se limita a ejecutar selectores:
se construye sobre datos validados, particiones sin fuga de información y
contrastes estadísticos explícitos, y mide propiedades que el rendimiento
no captura, como la estabilidad, la redundancia interna, la separación
frente al azar y la coherencia con las variables que usa el modelo. Sobre
esa base, los objetivos 4 a 6 implementan y evalúan el método cuántico
(su formulación, su codificación física y su contrato de evaluación), lo
comparan con los mejores selectores clásicos bajo el mismo protocolo y
añaden un control frente al óptimo exacto de la función de coste que
permite separar la calidad del criterio de la del optimizador. La
evaluación del método cuántico se realiza en un entorno simulado y bajo
una envolvente operativa controlada, por lo que las conclusiones se
refieren al protocolo experimental evaluado y no a una demostración en
hardware físico.

